\Titular*%
{miscelanea}%
{Música, frases e máis}%
{}%
{}%

\Opinion{Frases célebres}{Manuel Galán Rodríguez}

Inicio esta sección na revista Momentum achegando todas as frases que poida
entre a anterior edición da revista e a seguinte. Evidentemente, sempre se
admiten colaboracións para futuras edicións. Espero que vos guste!

\begin{multicols}{2}
\raggedcolumns

\Cita{¿Recordáis las palabras de Fernando Simón de que como mucho iba a haber 1
ó 2 casos? Pues en esta asignatura va a ser igual: va a haber como mucho 1 ó 2
suspensos.} {Diego Martínez,\\ Física computacional avanzada, 2025.}

% \Cita{Esta teoría una vez que uno la entiende, acaba reconfortado. Igual que
% uno acaba reconfortado después de haber hecho ejercicio físico.}
% {Juan Saborido,\\ Física de partículas I, 2025.}

% \Cita{Esto está explícitamente demostrado en muy pocos casos porque es
% horrible. [...] La vida es corta. Esto a los matemáticos les horroriza, pero
% después lo usan.} {Néstor,\\ Teoría cuántica de campos, 2025.}

\Cita{Igual estáis pensando que estoy haciendo una cosa extraña. Os lo va a
dejar de parecer dentro de un minuto porque la voy a complicar más.}
{Néstor,\\ Teoría cuántica de campos, 2025.}

% \Cita{Las matemáticas… las matemáticas son un lenguaje que de por sí solo no
% nos dice nada. Cuando queremos hacer poesía con ese lenguaje… lo llamamos
% Física.} {Luismi,\\ Mecánica estatística I, 2025.}

% \Cita{Que facemos? Seguimos ou xogamos unha partida de cartas?.}
% {Raúl,\\ Tecnoloxía do láser, 2025.}

\Cita{Y a mí [me da pereza] la física. Y la vida en general.}
{Víctor Pardo,\\ Electromagnetismo I, 2025.}

\Cita{El péndulo está un poco bailarín... no sabe para dónde va. Es como yo
hoy, no sé para dónde voy.}
{Jaime Álvarez,\\ Técnicas Experimentais II, 2025.}

\Cita{A potencia [solar] que chega aquí [polo norte] é moi pequena. Isto o
saben ben os pingüinos.}
{Carlos Montero,\\ Óptica I, 2025.}


\end{multicols}

\newpage
\Opinion{Praceres descoñecidos: $E=mc^2$ de B.A.D.}{Carlos Carballeira Romero}%

\begin{multicols}{2}

A expresión da equivalencia entre masa e enerxía é, posiblemente, a ecuación da
física máis implantada no imaxinario popular. Por iso, cando lle propuxen ao
equipo directivo de \textit{Momentum} realizar unha serie de artigos con
escolmas de cancións inspiradas na física, tiven bastante claro que $E=mc^2$ ía
ser o punto de partida. O título deste recuncho, \textit{Praceres
descoñecidos}, tampouco é casual... pero a súa explicación queda adiada para
unha futura entrega.

\begin{figure}[H]
    \centering
    \includegraphics[width=\linewidth]{revistas/004/imaxes/emc2.jpg}
    \caption*{
        Portada do \textit{single} $E=mc^2$ da banda \textit{B.A.D.} (\textit{Big
        Audio Dynamite}) que chegou ao posto 11 da \textit{UK singles Chart}. No
        centro (con gabardina) Mick Jones, membro fundador, cantante e compositor
        de \textit{The Clash}.
    }
\end{figure}

Froito da súa popularidade, $E=mc^2$ ten sido fonte de inspiración para
artistas/músicos de diferentes estilos que van dende o rap á música
electrónica, pasando polo pop de radiofórmula máis comercial encarnado por
Mariah Carey. Porque, por sorprendente que vos poida parecer, antes de se
converter na pregoeira oficial do nadal nos EEUU a cantante de Nova York
publicou non só unha canción, senón un álbum enteiro, titulado $E=mc^2$. Iso
si, Mariah deulle unha volta ás implicacións da celebérrima ecuación declarando
que o seu significado é \textit{emancipación igual a Mariah Carey multiplicado
por dous}. Queda claro que, antes de recibir clases de física, a boa de Mariah
ten que darlle un repaso xeral ás matemáticas básicas.

Descartada Mariah Carey pola súa heterodoxa reinterpretación da fórmula de
Einstein, había dúas serias candidaturas á primeira entrega desta sección. Unha
delas era o álbum $E=mc^2$ (con canción homónima) do italiano Giorgio Moroder,
todo un pioneiro da música electrónica, \textit{techno} e do \textit{synthpop}.
Porén, todo un gañador de tres premios \textit{Óscar} (un a mellor B.S.O., e
dous a mellor canción) quedou finalmente descartado pola miña escolla final: o
tema $E=mc^2$ que a banda \textit{B.A.D.} (acrónimo de \textit{Big Audio
Dynamite}) escolleu como segundo single do seu álbum de debut. O motivo desa
decisión é que \textit{B.A.D.} foi o proxecto que creou Mick Jones, membro
fundador, compositor e \textit{frontman} de \textit{The Clash}, logo de marchar
(ou de que o botasen) da mítica banda de punk-rock. E se $E=mc^2$ está
implantado no imaxinario popular, non o están menos os primeiros acordes de
\textit{Should I stay or should I go}, todo un himno creado por Mick Jones fai
máis de corenta anos que vai saltando de xeración en xeración. Esta conexión
entre un dos físicos máis influentes e unha das bandas máis influentes non se
podía deixar pasar por alto.

A saída de Jones de \textit{The Clash} veu acompañada tamén dunha muda cara a
sons máis comerciais (ou emerxentes) nos anos 80. De este xeito, sen abandonar
de todo o \textit{punk}, \textit{B.A.D.} experimentou co \textit{pop}, co
\textit{dance}, co \textit{hip hop} e co \textit{funk} cunha acollida favorable
por parte do público. Así, o \textit{single} $E=mc^2$ chegou ao número 11 na
lista de \textit{singles} do Reino Unido. Iso si, se botades un ollo ao vídeo
ou á letra da canción, non esperedes unha disertación sobre física: o tema está
inspirado en varias das películas do cineasta británico Nick Roeg e, de feito,
a maior parte da letra e do vídeo son, respectivamente, frases e secuencias
extraídas das mesmas. Unha delas, titulada \textit{Insignificance} (1985), ten
como argumento o hipotético encontro de catro iconas dos anos 50 do século
pasado nun hotel de Nova York: Marilyn Monroe, o seu marido e estrela do
béisbol Joe DiMaggio, o senador McCarthy (impulsor da caza de bruxas
anticomunista nos EEUU por aqueles anos) e Albert Einstein. Velaí por que a
canción se titula $E=mc^2$ e as alusións á relatividade e ao transcorrer do
tempo na letra. Alguén pensará que, entón, a relación coa física é feble. Se é
o voso caso, lembrade que, cando anos antes compuxo \textit{Should I stay or
should I go}, Mick Jones resumiu á perfección o dilema que afronta a diario
calquera postdoc ou estudante de doutoramento.

\end{multicols}
